%% Brochure marketing Kesh
%% Version canonique française
%% À mettre à jour à chaque release Kesh (cf. ../README.md)

\documentclass[a4paper,11pt,oneside]{article}

\usepackage[french]{babel}
\usepackage{../shared/kesh-style}
\usepackage{afterpage}
\usepackage{multicol}

\renewcommand{\Kesh@logo@path}{../shared/kesh-logo.svg}

% La brochure utilise des marges plus généreuses pour aérer
\geometry{
    a4paper,
    top=2.2cm,
    bottom=2.2cm,
    left=2.0cm,
    right=2.0cm,
    headheight=14pt,
    headsep=0.6cm,
    footskip=1cm
}

% Pas de TOC pour la brochure (format court et visuel)

% Pas de header/footer standards pour la brochure (look plus magazine)
\pagestyle{empty}

\begin{document}

%%─────────────────────────────────────────────────────────────────────────────
%% Page de couverture impactante full-bleed
%%─────────────────────────────────────────────────────────────────────────────

\begin{titlepage}
    \pagecolor{keshPrimary}\afterpage{\nopagecolor}
    \color{white}

    \vspace*{1.5cm}

    \begin{center}
        \begin{tcolorbox}[%
            enhanced, colback=white, colframe=white,
            width=5cm, height=5cm,
            boxsep=0pt, boxrule=0pt, arc=12pt,
            halign=center, valign=center%
        ]
            \textbf{\color{keshPrimary}\fontsize{60pt}{72pt}\selectfont\sffamily Kesh}
        \end{tcolorbox}

        \vspace{2.5cm}

        % Slogan principal
        {\sffamily\fontsize{40pt}{48pt}\selectfont\bfseries\color{white} Votre comptabilité\par
        \vspace{0.3cm}
        suisse, souveraine\par
        \vspace{0.3cm}
        et open-source.\par}

        \vspace{2cm}

        {\sffamily\Large\color{white}%
        Pour les PME, indépendants et associations\par
        en quête de simplicité, conformité et liberté.\par}

        \vfill

        \begin{tcolorbox}[%
            enhanced, colback=white!10!keshPrimary, colframe=white,
            width=12cm, boxsep=10pt, boxrule=1pt, arc=4pt%
        ]
            \begin{center}
                {\sffamily\large \textbf{Conformité} Code des obligations CH\par
                \vspace{4pt}
                \textbf{QR Bill} 2.2 \quad\textbar\quad
                \textbf{CAMT.053} \quad\textbar\quad
                \textbf{Multilingue}\par
                \vspace{4pt}
                \textbf{Self-hosting} Docker \quad\textbar\quad
                \textbf{Open-source} EUPL 1.2}
            \end{center}
        \end{tcolorbox}

        \vspace{0.5cm}

        {\sffamily\small\color{white}Version \textbf{\keshVersion} \quad\textbar\quad \keshReleaseDate \quad\textbar\quad \keshWebsite}
    \end{center}
\end{titlepage}

%%─────────────────────────────────────────────────────────────────────────────
%% Pourquoi Kesh — pitch principal
%%─────────────────────────────────────────────────────────────────────────────

\section*{\color{keshPrimary}Pourquoi Kesh ?}

Vous êtes \textbf{indépendant}, dirigeant de \textbf{PME suisse}, trésorier d'\textbf{association} ou fiduciaire gérant plusieurs dossiers clients ? Vous cherchez un logiciel de comptabilité qui soit à la fois :

\begin{multicols}{2}
\begin{itemize}
    \item[\faCheckCircle] Conforme au Code des obligations suisse
    \item[\faCheckCircle] Compatible QR Bill 2.2
    \item[\faCheckCircle] Capable d'importer du CAMT.053 ISO 20022
    \item[\faCheckCircle] Multilingue (FR, DE, IT, EN)
    \item[\faCheckCircle] Souverain (vos données restent chez vous)
    \item[\faCheckCircle] Open-source et auditable
    \item[\faCheckCircle] Sans frais récurrents par utilisateur
    \item[\faCheckCircle] Self-hostable avec Docker
\end{itemize}
\end{multicols}

\textbf{Kesh répond à tous ces critères.} Conçu spécifiquement pour le tissu PME suisse, ce logiciel combine la rigueur exigée par le droit comptable helvétique avec une expérience utilisateur moderne et fluide.

\vspace{0.5cm}

\begin{keshtip}
\textbf{Notre conviction.} La comptabilité d'une PME ne devrait pas dépendre d'un éditeur de logiciel basé à l'étranger, ni transmettre les données financières sensibles à un cloud propriétaire. Kesh est conçu pour respecter votre souveraineté : code source ouvert, déploiement chez vous, exports standards à tout moment.
\end{keshtip}

%%─────────────────────────────────────────────────────────────────────────────
%% Différenciateurs clés
%%─────────────────────────────────────────────────────────────────────────────

\clearpage

\section*{\color{keshPrimary}Ce qui distingue Kesh}

\vspace{0.5cm}

\begin{tcolorbox}[%
    enhanced, breakable,
    colback=keshSurfaceAlt, colframe=keshPrimary,
    boxrule=0pt, leftrule=4pt, arc=4pt,
    title=\sffamily\bfseries\color{keshPrimary}{\faUniversity\ Conformité réglementaire suisse}
]
\sffamily
Kesh implémente nativement les exigences du \textbf{Code des obligations suisse Art. 957a (tenue de la comptabilité) et Art. 958f (conservation 10 ans)}, ainsi que l'\textbf{Ordonnance OLICo RS 221.431} sur la conservation des livres de comptes.

L'\emph{audit-trail immutable} (insert-only) + le \emph{hash SHA-256} sur tous les exports garantissent l'intégrité légale au sens de l'OLICo Art. 9. Une recherche réglementaire dédiée (605 lignes documentées) accompagne le logiciel.
\end{tcolorbox}

\vspace{0.3cm}

\begin{tcolorbox}[%
    enhanced, breakable,
    colback=keshSurfaceAlt, colframe=keshSuccess,
    boxrule=0pt, leftrule=4pt, arc=4pt,
    title=\sffamily\bfseries\color{keshSuccess}{\faQrcode\ QR Bill 2.2 natif}
]
\sffamily
Génération automatique de PDF de factures conformes au \textbf{standard suisse QR Bill 2.2} (norme SIX 2020+). Compatible avec toutes les applications e-banking suisses : PostFinance, Raiffeisen, UBS, Credit Suisse, Migros Bank, Twint Pro.

Calcul automatique de la \emph{référence structurée} (ISO 11649 ou QRR selon votre IBAN), validation des données par checksum, support QR-IBAN et IBAN standard.
\end{tcolorbox}

\vspace{0.3cm}

\begin{tcolorbox}[%
    enhanced, breakable,
    colback=keshSurfaceAlt, colframe=keshInfo,
    boxrule=0pt, leftrule=4pt, arc=4pt,
    title=\sffamily\bfseries\color{keshInfo}{\faExchangeAlt\ Import bancaire automatisé}
]
\sffamily
Import natif des extraits de compte \textbf{CAMT.053 ISO 20022} (standard suisse officiel) et CSV multi-encodage avec \emph{profils banque réutilisables}. Détection automatique des doublons, gestion du rejet partiel.

Réconciliation \emph{automatique avec scoring} (date, montant, référence, libellé) et \emph{manuelle} pour les cas complexes. Éclatement de transactions agrégées (split). Règles d'affectation automatique pour les écritures récurrentes (loyers, salaires, abonnements).
\end{tcolorbox}

\vspace{0.3cm}

\begin{tcolorbox}[%
    enhanced, breakable,
    colback=keshSurfaceAlt, colframe=keshWarning,
    boxrule=0pt, leftrule=4pt, arc=4pt,
    title=\sffamily\bfseries\color{keshWarning}{\faLanguage\ Multilingue à 100\%}
]
\sffamily
Interface utilisateur entièrement traduite en \textbf{français, allemand, italien et anglais}. Idéal pour les PME couvrant plusieurs régions linguistiques suisses, les entreprises romandes en partenariat avec des fournisseurs alémaniques, ou les fiduciaires multilingues.

Les factures peuvent être émises dans la langue du client (différente de la langue d'interface). Les rapports comptables sont disponibles dans chacune des 4 langues.
\end{tcolorbox}

\vspace{0.3cm}

\begin{tcolorbox}[%
    enhanced, breakable,
    colback=keshSurfaceAlt, colframe=keshError,
    boxrule=0pt, leftrule=4pt, arc=4pt,
    title=\sffamily\bfseries\color{keshError}{\faLock\ Souveraineté des données}
]
\sffamily
Self-hosting via \textbf{Docker Compose} chez vous ou votre fiduciaire. Aucune donnée transmise à un cloud tiers sans votre consentement explicite.

Export ZIP de souveraineté en format ouvert (CSV + PDF + metadata SHA-256) téléchargeable à tout moment. Vous n'êtes \textbf{jamais prisonnier} de Kesh : migration vers un autre logiciel toujours possible.
\end{tcolorbox}

\vspace{0.3cm}

\begin{tcolorbox}[%
    enhanced, breakable,
    colback=keshSurfaceAlt, colframe=keshPrimaryLight,
    boxrule=0pt, leftrule=4pt, arc=4pt,
    title=\sffamily\bfseries\color{keshPrimaryLight}{\faUsers\ Multi-tenant pour fiduciaires}
]
\sffamily
Gérez \textbf{plusieurs sociétés} sur une seule instance Kesh. Isolation stricte des données par \texttt{company\_id} (défense en profondeur contre IDOR), audit-trail cross-company pour les super-administrateurs.

Idéal pour les \textbf{fiduciaires} souhaitant offrir Kesh comme solution à leurs clients tout en gardant une vue d'ensemble sur leur portefeuille de dossiers.
\end{tcolorbox}

%%─────────────────────────────────────────────────────────────────────────────
%% Cas d'usage par profil
%%─────────────────────────────────────────────────────────────────────────────

\clearpage

\section*{\color{keshPrimary}Pour qui est Kesh ?}

\vspace{0.5cm}

\begin{tcolorbox}[%
    enhanced, breakable,
    colback=white, colframe=keshBorder,
    boxrule=0.5pt, arc=4pt%
]
\sffamily

\textbf{\large\color{keshPrimary}{\faUserTie\ Indépendant ou entreprise individuelle}}

\vspace{0.3cm}

Vous êtes consultant, artisan, freelance, médecin, avocat, architecte ou autre indépendant suisse ? Kesh vous offre une solution comptable complète sans surdimensionnement :

\begin{itemize}
    \item Comptabilité simplifiée (recettes-dépenses) ou complète selon votre CA.
    \item Émission de factures QR Bill compatibles avec tous les e-banking.
    \item Suivi des paiements et relances clients facilités.
    \item Rapports de fin d'année pour votre fiduciaire ou votre déclaration personnelle.
\end{itemize}

\end{tcolorbox}

\vspace{0.5cm}

\begin{tcolorbox}[%
    enhanced, breakable,
    colback=white, colframe=keshBorder,
    boxrule=0.5pt, arc=4pt%
]
\sffamily

\textbf{\large\color{keshPrimary}{\faBuilding\ PME suisse (Sàrl, SA, coopérative)}}

\vspace{0.3cm}

Vous dirigez une PME avec 1--50 salariés et un chiffre d'affaires allant jusqu'à plusieurs millions ? Kesh vous accompagne dans :

\begin{itemize}
    \item Gestion de plan comptable conforme Sterchi PME ou KMU.
    \item Facturation client en lot avec QR Bill et suivi des échéances.
    \item Import bancaire et réconciliation en un clic.
    \item Production des comptes annuels (bilan, P\&L) conformes Art. 959/959a/959b CO.
    \item Multi-utilisateurs avec RBAC : comptable, opérateur de saisie, dirigeant en lecture seule.
\end{itemize}

\end{tcolorbox}

\vspace{0.5cm}

\begin{tcolorbox}[%
    enhanced, breakable,
    colback=white, colframe=keshBorder,
    boxrule=0.5pt, arc=4pt%
]
\sffamily

\textbf{\large\color{keshPrimary}{\faHandshake\ Association ou fondation}}

\vspace{0.3cm}

Trésorier d'une association sportive, culturelle, syndicale ou caritative ? Kesh gère parfaitement les structures à but non-lucratif :

\begin{itemize}
    \item Plan comptable adapté aux associations.
    \item Suivi des cotisations membres et dons.
    \item Comptes annuels lisibles pour l'assemblée générale.
    \item Conservation 10 ans conforme OLICo, même pour les petites associations dispensées d'audit (Art. 83b CC).
\end{itemize}

\end{tcolorbox}

\vspace{0.5cm}

\begin{tcolorbox}[%
    enhanced, breakable,
    colback=white, colframe=keshBorder,
    boxrule=0.5pt, arc=4pt%
]
\sffamily

\textbf{\large\color{keshPrimary}{\faBalanceScale\ Fiduciaire}}

\vspace{0.3cm}

Cabinet fiduciaire ou expert-comptable indépendant ? Kesh vous offre une plateforme moderne pour vos dossiers clients :

\begin{itemize}
    \item Mode multi-tenant : N companies sur une seule instance.
    \item Bascule rapide entre dossiers clients via le sélecteur de company.
    \item Isolation stricte des données (conformité protection des données et secret professionnel).
    \item Audit-trail centralisé pour la responsabilité fiduciaire.
    \item Tarification self-hosted : pas de coût par utilisateur ni par dossier.
\end{itemize}

\end{tcolorbox}

%%─────────────────────────────────────────────────────────────────────────────
%% Stack technique
%%─────────────────────────────────────────────────────────────────────────────

\clearpage

\section*{\color{keshPrimary}Une stack technique moderne et robuste}

\vspace{0.5cm}

Kesh repose sur des technologies éprouvées et modernes, choisies pour leur performance, leur sécurité et leur pérennité :

\vspace{0.5cm}

\begin{tabularx}{\linewidth}{>{\bfseries\color{keshPrimary}}lX}
    \toprule
    \keshtableheader{Couche} & \keshtableheader{Technologie} \\
    \midrule
    Backend & Rust 1.85+ (édition 2024) avec le framework Axum \\
    Frontend & SvelteKit 2 + Svelte 5, TypeScript, Tailwind CSS 4 \\
    Base de données & MariaDB 11.4 (compatibilité MySQL 8) \\
    Authentification & JWT + refresh tokens avec rotation \\
    Déploiement & Docker Compose, image officielle Docker Hub \\
    Internationalisation & Fluent (Mozilla) pour les traductions \\
    Génération PDF & Bibliothèque Rust native \\
    Tests & cargo test + Vitest + Playwright (E2E) \\
    CI/CD & GitHub Actions \\
    Licence & EUPL 1.2 (compatible GPL) \\
    \bottomrule
\end{tabularx}

\vspace{1cm}

\subsection*{\color{keshPrimary}Pourquoi ces choix ?}

\textbf{Rust} apporte la sécurité mémoire et la performance native (10--100× plus rapide que des stacks Node.js ou Python équivalentes). Pas de garbage collector, pas de race conditions silencieuses.

\textbf{Svelte 5} produit une interface réactive et rapide avec un bundle final compact (typiquement 50--150 Ko gzippé pour Kesh). Pas de virtual DOM coûteux comme React.

\textbf{MariaDB} offre une excellente performance pour la majorité des workloads PME, avec des optimisations spécifiques (FULLTEXT, JSON, REPEATABLE READ). Compatibilité MySQL 8 maintenue pour la portabilité.

\textbf{Docker Compose} simplifie le déploiement à un seul fichier \texttt{docker-compose.yml}. Idéal pour le self-hosting chez une PME ou un fiduciaire sans expertise Kubernetes.

%%─────────────────────────────────────────────────────────────────────────────
%% Roadmap visible
%%─────────────────────────────────────────────────────────────────────────────

\clearpage

\section*{\color{keshPrimary}Feuille de route}

Kesh est en développement actif. Voici la feuille de route publique au \keshReleaseDate :

\vspace{0.5cm}

\textbf{\color{keshSuccess}\faCheckCircle\ Déjà livré (Epic 1 à 9.5, v0.1 partielle)}

\begin{itemize}
    \item Fondations et authentification multi-tenant.
    \item Onboarding guidé et expert.
    \item Plan comptable suisse PME (Sterchi, KMU) avec saisie d'écritures partie double.
    \item Carnet d'adresses, catalogue produits, conditions de paiement.
    \item Facturation QR Bill 2.2 avec numérotation séquentielle automatique.
    \item Import bancaire CAMT.053 + CSV multi-encodage.
    \item Réconciliation automatique avec scoring + manuelle + règles d'affectation.
    \item Éclatement de transactions agrégées.
    \item Rapports comptables (Bilan, Compte de résultat, Balance, Journal).
    \item Exports PDF et CSV par rapport + ZIP global de souveraineté avec SHA-256.
\end{itemize}

\vspace{0.4cm}

\textbf{\color{keshInfo}\faCog\ En préparation (Epic 10, prochaine étape)}

\begin{itemize}
    \item Déploiement \& Opérations : CI/CD release automatisée, images Docker officielles, migrations automatiques, résilience frontend, guide d'installation PDF.
\end{itemize}

\vspace{0.4cm}

\textbf{\color{keshWarning}\faClock\ Backlog v0.1 (Epic 11 à 15)}

\begin{itemize}
    \item TVA Suisse : configuration des taux + calcul + rapport par période AFC.
    \item Avoirs et notes de crédit dédiés + fichiers de paiement pain.001.001.03 (ISO 20022).
    \item Budgets annuels + suivi.
    \item Clôture d'exercice avec report de soldes automatique.
    \item Justificatifs, lettrage, journaux personnalisables.
\end{itemize}

\vspace{0.4cm}

\textbf{\color{keshPrimaryLight}\faRocket\ Vision v0.2}

\begin{itemize}
    \item Signature électronique qualifiée (QES) sur les exports.
    \item API externes avec clés PAT pour intégrations IA et logiciels tiers.
    \item Multi-devises pour les PME à l'export.
    \item Format PDF/A (ISO 19005) pour l'archivage durable.
    \item Traductions DE/IT/EN des manuels.
    \item Webhooks et MCP server pour assistants IA conversationnels.
\end{itemize}

\vspace{0.5cm}

\begin{keshnote}
Notre feuille de route est publique sur GitHub. Suivez notre progression sur \keshWebsite ou consultez la documentation projet pour le détail des epics et stories.
\end{keshnote}

%%─────────────────────────────────────────────────────────────────────────────
%% Tarification et licence
%%─────────────────────────────────────────────────────────────────────────────

\clearpage

\section*{\color{keshPrimary}Tarification et licence}

\vspace{0.5cm}

\begin{tcolorbox}[%
    enhanced,
    colback=keshSurfaceAlt, colframe=keshPrimary,
    boxrule=2pt, arc=8pt,
    halign=center%
]
\sffamily

\vspace{0.3cm}

{\Huge\bfseries\color{keshPrimary}Kesh est gratuit.}

\vspace{0.4cm}

{\large Aucun coût de licence, aucun abonnement, aucun frais par utilisateur.}

\vspace{0.4cm}

{\normalsize Distribué sous licence \textbf{EUPL 1.2} (European Union Public Licence).\\
Vous pouvez l'utiliser, le modifier et le redistribuer librement, y compris en contexte commercial.}

\vspace{0.3cm}

\end{tcolorbox}

\vspace{0.5cm}

\textbf{Coûts éventuels associés au déploiement} (totalement indépendants de Kesh) :

\begin{itemize}
    \item Un serveur ou VPS pour héberger Kesh (~5--20 CHF/mois selon hébergeur).
    \item Un nom de domaine et un certificat HTTPS (Let's Encrypt gratuit).
    \item Optionnel : services de sauvegarde cloud externe (~5--15 CHF/mois selon volume).
    \item Optionnel : prestation d'installation et maintenance par un fiduciaire / ESN partenaire.
\end{itemize}

\vspace{0.5cm}

\subsection*{\color{keshPrimary}Pourquoi gratuit ?}

Kesh est conçu comme un \emph{commun numérique} pour le tissu économique suisse. Notre conviction : la comptabilité d'une PME relève de l'infrastructure essentielle, au même titre que la route ou le réseau électrique. Elle ne devrait pas dépendre d'une rente perpétuelle versée à un éditeur étranger.

Le développement de Kesh est financé par contributions volontaires, prestations partenaires (installation, formation, support) et soutiens de la communauté open-source. Si vous souhaitez contribuer financièrement ou techniquement, contactez l'équipe via GitHub.

%%─────────────────────────────────────────────────────────────────────────────
%% Call-to-action et contact
%%─────────────────────────────────────────────────────────────────────────────

\clearpage

\section*{\color{keshPrimary}Commencer avec Kesh}

\vspace{1cm}

\begin{tcolorbox}[%
    enhanced,
    colback=keshPrimary, colframe=keshPrimary,
    boxrule=0pt, arc=8pt,
    halign=center,
    coltext=white,
    boxsep=20pt%
]
\sffamily

{\Huge\bfseries\color{white}3 façons de découvrir Kesh}

\end{tcolorbox}

\vspace{0.8cm}

\begin{tabularx}{\linewidth}{cX}
    \begin{minipage}{2.8cm}
        \begin{tcolorbox}[%
            enhanced, colback=keshPrimary, colframe=keshPrimary,
            width=2.5cm, height=2.5cm,
            halign=center, valign=center,
            boxsep=0pt, boxrule=0pt, arc=10pt,
            coltext=white%
        ]
            {\fontsize{40pt}{40pt}\selectfont\faGithub}
        \end{tcolorbox}
    \end{minipage}
    &
    \begin{minipage}{12cm}
        \textbf{\large\color{keshPrimary}1. Explorez le code source}

        \vspace{0.3cm}

        Consultez le dépôt GitHub officiel à \href{\keshWebsite}{\keshWebsite}. Vous y trouverez le code source complet, la documentation technique, la feuille de route détaillée, les issues ouvertes et l'historique de développement.

        Vous pouvez \texttt{git clone} le dépôt et lancer une instance locale en moins de 5 minutes avec \texttt{docker compose up}.
    \end{minipage}
\end{tabularx}

\vspace{1cm}

\begin{tabularx}{\linewidth}{cX}
    \begin{minipage}{2.8cm}
        \begin{tcolorbox}[%
            enhanced, colback=keshSuccess, colframe=keshSuccess,
            width=2.5cm, height=2.5cm,
            halign=center, valign=center,
            boxsep=0pt, boxrule=0pt, arc=10pt,
            coltext=white%
        ]
            {\fontsize{40pt}{40pt}\selectfont\faRocket}
        \end{tcolorbox}
    \end{minipage}
    &
    \begin{minipage}{12cm}
        \textbf{\large\color{keshPrimary}2. Déployez votre instance}

        \vspace{0.3cm}

        Suivez le \emph{Manuel administrateur Kesh} pour installer une instance complète sur votre serveur ou VPS. Configuration en moins d'une heure si vous êtes familier avec Docker.

        Pour une assistance professionnelle (installation, migration depuis un autre logiciel, formation), contactez un \emph{partenaire Kesh certifié} via notre annuaire en ligne.
    \end{minipage}
\end{tabularx}

\vspace{1cm}

\begin{tabularx}{\linewidth}{cX}
    \begin{minipage}{2.8cm}
        \begin{tcolorbox}[%
            enhanced, colback=keshInfo, colframe=keshInfo,
            width=2.5cm, height=2.5cm,
            halign=center, valign=center,
            boxsep=0pt, boxrule=0pt, arc=10pt,
            coltext=white%
        ]
            {\fontsize{40pt}{40pt}\selectfont\faComments}
        \end{tcolorbox}
    \end{minipage}
    &
    \begin{minipage}{12cm}
        \textbf{\large\color{keshPrimary}3. Rejoignez la communauté}

        \vspace{0.3cm}

        Posez vos questions, partagez vos retours d'expérience, contribuez au code ou à la documentation : la communauté Kesh vous accueille sur \href{\keshWebsite/discussions}{GitHub Discussions} et \href{\keshWebsite/issues}{GitHub Issues}.

        Toute contribution est bienvenue : rapport de bug, demande de fonctionnalité, traduction, amélioration de la documentation, ou pull request de code.
    \end{minipage}
\end{tabularx}

\vfill

%%─────────────────────────────────────────────────────────────────────────────
%% Footer de la dernière page
%%─────────────────────────────────────────────────────────────────────────────

\begin{tcolorbox}[%
    enhanced,
    colback=keshSurfaceAlt, colframe=keshBorder,
    boxrule=0.5pt, arc=4pt,
    halign=center%
]
\sffamily\small\color{keshTextMuted}

\vspace{0.2cm}

\textbf{\color{keshPrimary}\Large Kesh}

\vspace{0.2cm}

Logiciel libre de comptabilité PME suisse \quad\textbar\quad Licence \href{https://joinup.ec.europa.eu/collection/eupl/eupl-text-eupl-12}{EUPL 1.2}

\vspace{0.2cm}

Version \textbf{\keshVersion} \quad\textbar\quad Publié le \keshReleaseDate \quad\textbar\quad \keshWebsite

\vspace{0.3cm}

\emph{Cette brochure marketing décrit Kesh à la date de publication ci-dessus. Pour la dernière version, consultez le site officiel. La brochure est mise à jour à chaque release.}

\vspace{0.2cm}

\end{tcolorbox}

\end{document}
